\documentclass[../main.tex]{subfiles}

\begin{document}

\chapter{Probability}

\begin{definition}{Discrete Probability Model}{}
    Let \(   \Omega   \) be a set , \(  \left|  \Omega  \right|< \infty   \) .  \(  Q=  \varphi \left(  \Omega  \right)   \) the collection of all subsets of \(   \Omega   \) .  \(  P:  \varphi \left( E \right)\to \left[ 0,1 \right]    \)       a function . If
    \begin{enumerate}
        \item \(  P\left(  \Omega  \right)= 1   \).
        \item  \(  P\left( A\cup B \right)= P\left( A \right)+ P\left( B \right)     \) for all \(  A,B \in  \varphi \left(  \Omega  \right)   \) , \(  A\cap B= \varnothing  \) .     
    \end{enumerate}
    Then we call \(  \left(  \Omega ,Q,P \right)   \) a  discrete probability model . 
\end{definition}


\begin{proposition}{}{}
    Let \(  \left(  \Omega ,Q,P \right)   \) be a discerte probability model , 
    \begin{enumerate}
        \item Let \(  A \in  \varphi \left(  \Omega  \right)    \) ,  \(  \bar{A}: = \left\{ x \in  \Omega ,x \not\in A \right\}  \). THen \[
        P\left( \bar{A} \right)= 1-P\left( A \right)  
        \]
        \item \(  A,B \in  \varphi \left(  \Omega  \right)   \) , then \[
        P\left( A\cup B \right)= P\left( A \right)+ P\left( B \right)-P\left( A\cap B \right)    
        \]   
    \end{enumerate}
     
\end{proposition}

\begin{remark}
    If \(  P\left( \left\{ \cdot  \right\} \right)=  \frac{1 }{\left|  \Omega  \right|  }    \) , then we say the probability is \dfntxt{uniform} . 
\end{remark}

\begin{example}{}{}
    Take \(   \Omega   = \left\{ a,b \right\}\)  , \( Q=  P\left(  \Omega  \right)= \left\{ \varnothing, a,b,  \Omega  \right\}   \). \(  P\left( a \right)=  \theta , P\left( b \right)= 1- \theta     \) .  Then \(  P  \) can extends to a probability, and \(  \left(  \Omega ,Q,P  \right)   \) comes to a discerte probability model.     
\end{example}


\begin{definition}{Random Variable}{}
    A \dfntxt{Random Variable} (R.V) is  just a  map \(  X: \Omega \to \mathbb{R}   \) .  
\end{definition}

\begin{example}{A Random Variable}{}
    \(  A \in  \varphi \left(  \Omega  \right)   \) . Let \(  X =  \chi _{A}  \) . Then \(  X  \) is a random variable. And \[
    P\left( X= 1 \right)= P\left( A \right),\quad P\left( X= 0 \right)= 1-P\left( A \right).     
    \]
\end{example}

\begin{definition}{Expectation (Mean Value of a R.V)}{}
    Let \(  X  \) be a \(  R. V  \) . If 
    \begin{enumerate}
        \item \(  \left| X\left(  \Omega  \right)  \right|< \infty   \)  and \(  X\ge 0  \) .  We define  \[
        E\left( X \right): = \sum _{x \in X\left(  \Omega  \right) }x P\left( X= x \right)   
        \]
        \item In general case, if \(  E\left( \left| X \right|  \right)< \infty   \) , we define \[
        E\left( X \right): = E\left( X^{+ }\right)-E\left( X^{-}\right)   
        \] 
    \end{enumerate}
      
\end{definition}

\begin{proposition}{}{}
    Let \(   \varphi : X\left(  \Omega  \right)\to \mathbb{R}    \) . Assume that \(  E\left( \left|  \varphi \left( X \right) \right|  \right)< \infty   \). Then , \[
    E\left(  \varphi \left( x \right)  \right) =  \sum _{x \in X\left(  \Omega  \right) } \varphi \left( x \right)P\left( X= x \right)  
    \] 
\end{proposition}

\begin{definition}{}{}
    Suppose that \(  E\left( X^{2} \right)< \infty   \) . Then define the \dfntxt{variance} of \(  X  \)  \[
 \operatorname{Var}\left( X \right)  := E\left[ (X - E[X])^2 \right]= E\left[ X^{2} \right]-\left( E\left[ X \right]  \right)^{2}  
    \]  
\end{definition}
\begin{proof}
    \[
    \begin{aligned}
    \operatorname{Var}X= E\left[ \left( X-E\left[ X \right]  \right)^{2}  \right]&=  E\left[ X^{2}- 2E\left( X \right)X+ \left( E\left( X \right)  \right)^{2}   \right]\\ 
     &=  E\left( X^{2} \right)+ E\left( -2E\left( X \right)X  \right)+ E\left( \left( E\left[ X \right]  \right)^{2}  \right)     
    \end{aligned}
    \]
\end{proof}

\begin{definition}{Probability Generating Function}{}
    Let $X$ be a discrete random variable with range in $\{0, 1, 2, \dots\}$. Its Probability Generating Function, $G_X(t)$, is defined as:
$$G_X(t) := E[t^X]$$
for all real numbers $t$ for which the expectation exists. Specifically, 
$$G_X(t) = \sum_{k=0}^{\infty} t^k P(X=k)$$
\end{definition}

\begin{proposition}{}{}
    Assume that \(  R> 1  \) , Then \begin{enumerate}
        \item \(  E\left( X \right)= G_{X}^{\prime} \left( 1 \right)    \)
        \item  $\text{Var}(X) = G_X''(1) + G_X'(1) - [G_X'(1)]^2$
    \end{enumerate}
     
\end{proposition}

\begin{example}{Poisson}{}
    Let  \(  X\sim \operatorname{Poisson}\left(  \lambda  \right)   \) , $$P(X=k) = \frac{\lambda^k e^{-\lambda}}{k!}, \quad \text{for } k = 0, 1, 2, \dots$$
    \[
\begin{aligned} G_X(t) &= E[t^X] = \sum_{k=0}^{\infty} t^k P(X=k) \\ &= \sum_{k=0}^{\infty} t^k \frac{\lambda^k e^{-\lambda}}{k!} \\ &= e^{-\lambda} \sum_{k=0}^{\infty} \frac{(t\lambda)^k}{k!} \quad \text{(将常数和 t 移入)} \end{aligned}
\] \[
G_X(t) = e^{-\lambda} \cdot e^{t\lambda} = e^{t\lambda - \lambda} = e^{\lambda(t-1)}
\] \[
G_X(t) = e^{\lambda(t-1)}
\]
\end{example}

\end{document}